% ============================================================
%  Part XVII — Vocabulary
%
%  Nouns are written with their article and coloured by gender:
%    \dm{}  der (blue)   \df{}  die (red)
%    \dn{}  das (green)  \dpl{} die, plural only (amber)
%  Verbs: \vr{} regular (plain), \vi{} irregular (bold, accent colour).
% ============================================================
\gpart{XVII}{Vocabulary}

% ------------------------------------------------------------
\gsec{People and family}

\gintro{Every noun is given with its article, and the article is colour-coded:
\dm{Mann} blue, \df{Frau} red, \dn{Kind} green, \dpl{Leute} amber for nouns that
exist only in the plural.}

\begin{dtbl}[\footnotesize]{W{1.08}W{0.92}W{1.08}W{0.92}}
\rowcolor{Ink} \hd{Deutsch} & \hd{English} & \hd{Deutsch} & \hd{English} \\ \hdrule
\dm{Mann}        & man          & \df{Frau}         & woman \\
\dn{Kind}        & child        & \dpl{Leute}       & people \\
\dm{Junge}       & boy          & \dn{Mädchen}      & girl \\
\dm{Vater}       & father       & \df{Mutter}       & mother \\
\dpl{Eltern}     & parents      & \dpl{Geschwister} & siblings \\
\dm{Bruder}      & brother      & \df{Schwester}    & sister \\
\dm{Sohn}        & son          & \df{Tochter}      & daughter \\
\dm{Großvater}   & grandfather  & \df{Großmutter}   & grandmother \\
\dpl{Großeltern} & grandparents & \dn{Enkelkind}    & grandchild \\
\dm{Onkel}       & uncle        & \df{Tante}        & aunt \\
\dm{Cousin}      & cousin \eng{m} & \df{Cousine}    & cousin \eng{f} \\
\dm{Neffe}       & nephew       & \df{Nichte}       & niece \\
\dm{Ehemann}     & husband      & \df{Ehefrau}      & wife \\
\dm{Freund}      & friend \eng{m} & \df{Freundin}   & friend \eng{f} \\
\dm{Nachbar}     & neighbour    & \df{Familie}      & family \\
\dm{Mensch}      & human being  & \df{Person}       & person \\
\dm{Erwachsene}  & adult        & \dn{Baby}         & baby \\
\dm{Jugendliche} & teenager     & \dn{Ehepaar}      & married couple \\
\dm{Gast}        & guest        & \df{Hochzeit}     & wedding \\
\dm{Name}        & name         & \dn{Alter}        & age \\
\dm{Geburtstag}  & birthday     & \df{Geburt}       & birth \\
\dm{Chef}        & boss         & \df{Kollegin}     & colleague \eng{f} \\
\dm{Herr}        & Mr, gentleman & \df{Dame}        & lady \\
\dm{Partner}     & partner      & \df{Beziehung}    & relationship \\
\dn{Paar}        & couple, pair & \df{Liebe}        & love \\
\dm{Schwiegervater} & father-in-law & \df{Schwiegermutter} & mother-in-law \\
\dm{Nachname}    & surname      & \dm{Vorname}      & first name \\
\end{dtbl}

\begin{gnote}
\textit{Geschwister} is normally used in the plural (\textit{meine
Geschwister}). \textit{der Erwachsene}, \textit{der Jugendliche} and
\textit{der Verwandte} are adjectival nouns, so they decline like adjectives:
\textit{ein Erwachsen\textbf{er}}.
\end{gnote}

% ------------------------------------------------------------
\gsec{The body}

\begin{dtbl}[\footnotesize]{W{1.08}W{0.92}W{1.08}W{0.92}}
\rowcolor{Ink} \hd{Deutsch} & \hd{English} & \hd{Deutsch} & \hd{English} \\ \hdrule
\dm{Kopf}      & head       & \dn{Gesicht}  & face \\
\dn{Haar}      & hair       & \df{Stirn}    & forehead \\
\dn{Auge}      & eye        & \dn{Ohr}      & ear \\
\df{Nase}      & nose       & \dm{Mund}     & mouth \\
\df{Lippe}     & lip        & \dm{Zahn}     & tooth \\
\df{Zunge}     & tongue     & \dm{Hals}     & neck, throat \\
\df{Schulter}  & shoulder   & \dm{Arm}      & arm \\
\dm{Ellbogen}  & elbow      & \df{Hand}     & hand \\
\dm{Finger}    & finger     & \dm{Daumen}   & thumb \\
\df{Brust}     & chest      & \dm{Rücken}   & back \\
\dm{Bauch}     & stomach    & \df{Taille}   & waist \\
\dn{Bein}      & leg        & \dn{Knie}     & knee \\
\dm{Fuß}       & foot       & \df{Zehe}     & toe \\
\df{Haut}      & skin       & \dm{Knochen}  & bone \\
\dn{Herz}      & heart      & \df{Lunge}    & lung \\
\dm{Magen}     & stomach \eng{organ} & \df{Leber} & liver \\
\dn{Blut}      & blood      & \dm{Muskel}   & muscle \\
\dn{Gehirn}    & brain      & \dm{Nerv}     & nerve \\
\dm{Körper}    & body       & \df{Figur}    & figure \\
\dm{Bart}      & beard      & \df{Wange}    & cheek \\
\dn{Kinn}      & chin       & \df{Augenbraue} & eyebrow \\
\dm{Nagel}     & nail       & \df{Faust}    & fist \\
\dn{Handgelenk} & wrist     & \df{Hüfte}     & hip \\
\df{Ferse}     & heel       & \dm{Knöchel}  & ankle \\
\df{Stimme}    & voice      & \dm{Atem}     & breath \\
\dm{Schmerz}   & pain       & \df{Träne}    & tear \\
\dm{Schlaf}    & sleep      & \dm{Traum}    & dream \\
\end{dtbl}

\begin{gnote}
Body parts usually take the definite article where English uses a possessive:
\textit{Ich wasche mir \textbf{die} Hände} — \emph{I wash \textbf{my} hands}.
The owner appears as a dative pronoun instead.
\end{gnote}

% ------------------------------------------------------------
\gsec{Food and drink}

\begin{dtbl}[\footnotesize]{W{1.08}W{0.92}W{1.08}W{0.92}}
\rowcolor{Ink} \hd{Deutsch} & \hd{English} & \hd{Deutsch} & \hd{English} \\ \hdrule
\dn{Brot}      & bread      & \dn{Brötchen} & bread roll \\
\df{Butter}    & butter     & \dm{Käse}     & cheese \\
\dn{Ei}        & egg        & \df{Milch}    & milk \\
\dn{Fleisch}   & meat       & \dm{Fisch}    & fish \\
\dn{Hähnchen}  & chicken    & \dn{Schweinefleisch} & pork \\
\dn{Rindfleisch} & beef     & \df{Wurst}    & sausage \\
\dm{Schinken}  & ham        & \dm{Speck}    & bacon \\
\dn{Obst}      & fruit      & \dn{Gemüse}   & vegetables \\
\dm{Apfel}     & apple      & \df{Banane}   & banana \\
\df{Orange}    & orange     & \df{Zitrone}  & lemon \\
\df{Erdbeere}  & strawberry & \df{Traube}   & grape \\
\df{Kartoffel} & potato     & \df{Tomate}   & tomato \\
\df{Zwiebel}   & onion      & \dm{Knoblauch} & garlic \\
\dm{Salat}     & salad      & \df{Gurke}    & cucumber \\
\df{Karotte}   & carrot     & \dm{Pilz}     & mushroom \\
\dm{Reis}      & rice       & \df{Nudel}    & noodle, pasta \\
\dn{Mehl}      & flour      & \dm{Zucker}   & sugar \\
\dn{Salz}      & salt       & \dm{Pfeffer}  & pepper \\
\dn{Öl}        & oil        & \dm{Essig}    & vinegar \\
\df{Suppe}     & soup       & \df{Soße}     & sauce \\
\dm{Kuchen}    & cake       & \df{Schokolade} & chocolate \\
\dn{Eis}       & ice cream  & \dm{Keks}     & biscuit \\
\dn{Wasser}    & water      & \dm{Kaffee}   & coffee \\
\dm{Tee}       & tea        & \dm{Saft}     & juice \\
\dn{Bier}      & beer       & \dm{Wein}     & wine \\
\dn{Frühstück} & breakfast  & \dn{Mittagessen} & lunch \\
\dn{Abendessen} & dinner    & \df{Mahlzeit} & meal \\
\dm{Teller}    & plate      & \dn{Glas}     & glass \\
\df{Gabel}     & fork       & \dn{Messer}   & knife \\
\end{dtbl}

\begin{gnote}
\textit{Obst} and \textit{Gemüse} are uncountable and stay singular:
\textit{Ich esse viel Obst}. To count them you need a piece word, such as
\textit{ein Stück Kuchen} or \textit{eine Scheibe Brot}.
\end{gnote}

% ------------------------------------------------------------
\gsec{Animals}

\begin{dtbl}[\footnotesize]{W{1.08}W{0.92}W{1.08}W{0.92}}
\rowcolor{Ink} \hd{Deutsch} & \hd{English} & \hd{Deutsch} & \hd{English} \\ \hdrule
\dn{Tier}      & animal     & \dn{Haustier} & pet \\
\dm{Hund}      & dog        & \df{Katze}    & cat \\
\dn{Pferd}     & horse      & \df{Kuh}      & cow \\
\dn{Schwein}   & pig        & \dn{Schaf}    & sheep \\
\df{Ziege}     & goat       & \dn{Huhn}     & chicken \\
\dm{Hahn}      & cockerel   & \df{Ente}     & duck \\
\df{Gans}      & goose      & \dm{Vogel}    & bird \\
\df{Maus}      & mouse      & \df{Ratte}    & rat \\
\dn{Kaninchen} & rabbit     & \dm{Hase}     & hare \\
\dm{Fuchs}     & fox        & \dm{Wolf}     & wolf \\
\dm{Bär}       & bear       & \dm{Hirsch}   & deer \\
\dn{Reh}       & roe deer   & \dn{Eichhörnchen} & squirrel \\
\dm{Igel}      & hedgehog   & \dm{Maulwurf} & mole \\
\df{Fledermaus} & bat       & \df{Eule}     & owl \\
\dm{Adler}     & eagle      & \df{Taube}    & pigeon \\
\dm{Spatz}     & sparrow    & \df{Krähe}    & crow \\
\dm{Fisch}     & fish       & \dm{Hai}      & shark \\
\dm{Wal}       & whale      & \dm{Delfin}   & dolphin \\
\dm{Frosch}    & frog       & \df{Kröte}    & toad \\
\df{Schlange}  & snake      & \df{Eidechse} & lizard \\
\df{Schildkröte} & tortoise & \df{Spinne}   & spider \\
\df{Biene}     & bee        & \df{Wespe}    & wasp \\
\df{Ameise}    & ant        & \df{Fliege}   & fly \\
\df{Mücke}     & mosquito   & \dm{Schmetterling} & butterfly \\
\dm{Käfer}     & beetle     & \dm{Wurm}     & worm \\
\dm{Löwe}      & lion       & \dm{Tiger}    & tiger \\
\dm{Elefant}   & elephant   & \df{Giraffe}  & giraffe \\
\dm{Affe}      & monkey     & \dn{Zebra}    & zebra \\
\dn{Nilpferd}  & hippo      & \dn{Krokodil} & crocodile \\
\end{dtbl}

\begin{gnote}
Several animal names are weak nouns and add \textit{-n} or \textit{-en} in every
case but the nominative singular: \textit{der Hase}, \textit{der Löwe},
\textit{der Affe}, \textit{der Bär}, \textit{der Elefant}. \textit{Ich sehe den
Löw\textbf{en}}.
\end{gnote}

% ------------------------------------------------------------
\gsec{The house}

\begin{dtbl}[\footnotesize]{W{1.08}W{0.92}W{1.08}W{0.92}}
\rowcolor{Ink} \hd{Deutsch} & \hd{English} & \hd{Deutsch} & \hd{English} \\ \hdrule
\dn{Haus}       & house       & \df{Wohnung}   & flat \\
\dn{Zimmer}     & room        & \df{Küche}     & kitchen \\
\dn{Bad}        & bathroom    & \dn{Schlafzimmer} & bedroom \\
\dn{Wohnzimmer} & living room & \dm{Flur}      & hall \\
\dm{Keller}     & cellar      & \dm{Dachboden} & attic \\
\df{Garage}     & garage      & \dm{Garten}    & garden \\
\dm{Balkon}     & balcony     & \df{Terrasse}  & terrace \\
\df{Tür}        & door        & \dn{Fenster}   & window \\
\df{Wand}       & wall        & \dm{Boden}     & floor \\
\df{Decke}      & ceiling     & \dn{Dach}      & roof \\
\df{Treppe}     & staircase   & \dm{Aufzug}    & lift \\
\dm{Tisch}      & table       & \dm{Stuhl}     & chair \\
\dn{Sofa}       & sofa        & \dm{Sessel}    & armchair \\
\dn{Bett}       & bed         & \dm{Schrank}   & cupboard \\
\dn{Regal}      & shelf       & \df{Lampe}     & lamp \\
\dm{Teppich}    & carpet      & \dm{Vorhang}   & curtain \\
\dm{Spiegel}    & mirror      & \dn{Bild}      & picture \\
\df{Uhr}        & clock       & \dm{Schlüssel} & key \\
\dm{Kühlschrank} & fridge     & \dm{Herd}      & cooker \\
\dm{Ofen}       & oven        & \df{Spülmaschine} & dishwasher \\
\df{Waschmaschine} & washing machine & \dn{Waschbecken} & sink \\
\df{Dusche}     & shower      & \df{Badewanne} & bathtub \\
\df{Toilette}   & toilet      & \dn{Handtuch}  & towel \\
\df{Seife}      & soap        & \dn{Kissen}    & cushion \\
\df{Bettdecke}  & duvet       & \dn{Bettlaken} & sheet \\
\df{Heizung}    & heating     & \df{Steckdose} & socket \\
\dm{Müll}       & rubbish     & \df{Miete}     & rent \\
\dm{Schreibtisch} & desk      & \dm{Vermieter} & landlord \\
\end{dtbl}

\begin{gnote}
\textit{die Decke} means both \emph{ceiling} and \emph{blanket}; the bedding
sense is usually made explicit as \textit{die Bettdecke}. Note also
\textit{das Zimmer} (a room in general) beside \textit{der Raum} (a space).
\end{gnote}

% ------------------------------------------------------------
\gsec{Clothing}

\begin{dtbl}[\footnotesize]{W{1.08}W{0.92}W{1.08}W{0.92}}
\rowcolor{Ink} \hd{Deutsch} & \hd{English} & \hd{Deutsch} & \hd{English} \\ \hdrule
\df{Kleidung}   & clothing    & \dn{Kleid}     & dress \\
\dn{Hemd}       & shirt       & \df{Bluse}     & blouse \\
\df{Hose}       & trousers    & \dm{Rock}      & skirt \\
\dm{Pullover}   & jumper      & \dn{T-Shirt}   & T-shirt \\
\df{Jacke}      & jacket      & \dm{Mantel}    & coat \\
\dm{Anzug}      & suit        & \df{Krawatte}  & tie \\
\dpl{Jeans}     & jeans       & \dpl{Shorts}   & shorts \\
\dm{Schal}      & scarf       & \df{Mütze}     & woolly hat \\
\dm{Hut}        & hat         & \dm{Handschuh} & glove \\
\df{Socke}      & sock        & \dm{Schuh}     & shoe \\
\dm{Stiefel}    & boot        & \df{Sandale}   & sandal \\
\dm{Gürtel}     & belt        & \df{Tasche}    & bag \\
\dm{Rucksack}   & rucksack    & \df{Brille}    & glasses \\
\dm{Schmuck}    & jewellery   & \dm{Ring}      & ring \\
\df{Kette}      & necklace    & \df{Armbanduhr} & watch \\
\dm{Knopf}      & button      & \dm{Reißverschluss} & zip \\
\df{Größe}      & size        & \dm{Ärmel}     & sleeve \\
\dm{Kragen}     & collar      & \df{Naht}      & seam \\
\dm{Stoff}      & fabric      & \df{Baumwolle} & cotton \\
\df{Wolle}      & wool        & \dn{Leder}     & leather \\
\df{Seide}      & silk        & \df{Wäsche}    & laundry \\
\dm{Schlafanzug} & pyjamas    & \dm{Badeanzug} & swimsuit \\
\df{Unterwäsche} & underwear  & \df{Strumpfhose} & tights \\
\df{Umkleidekabine} & fitting room & \dm{Kleiderbügel} & coat hanger \\
\df{Mode}       & fashion     & \dm{Schneider} & tailor \\
\df{Wäscherei}  & laundry     & \dm{Fleck}     & stain \\
\df{Falte}      & crease      & \dn{Muster}    & pattern \\
\dm{Absatz}     & heel        & \dm{Schnürsenkel} & shoelace \\
\end{dtbl}

\begin{gnote}
\textit{die Hose} is singular in German although \emph{trousers} is plural in
English: \textit{Die Hose \textbf{ist} neu}. \textit{die Brille} behaves the
same way. \textit{die Jeans} and \textit{die Shorts} are borrowed with their
English plural.
\end{gnote}

% ------------------------------------------------------------
\gsec{Jobs and work}

\begin{dtbl}[\footnotesize]{W{1.08}W{0.92}W{1.08}W{0.92}}
\rowcolor{Ink} \hd{Deutsch} & \hd{English} & \hd{Deutsch} & \hd{English} \\ \hdrule
\dm{Beruf}      & profession  & \df{Arbeit}    & work \\
\dm{Arzt}       & doctor      & \df{Ärztin}    & doctor \eng{f} \\
\dm{Lehrer}     & teacher     & \df{Lehrerin}  & teacher \eng{f} \\
\dm{Ingenieur}  & engineer    & \dm{Architekt} & architect \\
\dm{Anwalt}     & lawyer      & \dm{Richter}   & judge \\
\df{Krankenschwester} & nurse \eng{f} & \dm{Krankenpfleger} & nurse \eng{m} \\
\dm{Polizist}   & police officer & \dm{Feuerwehrmann} & firefighter \\
\dm{Koch}       & cook        & \dm{Kellner}   & waiter \\
\dm{Bäcker}     & baker       & \dm{Metzger}   & butcher \\
\dm{Verkäufer}  & shop assistant & \dm{Kaufmann} & businessman \\
\dm{Bauer}      & farmer      & \dm{Gärtner}   & gardener \\
\dm{Tischler}   & carpenter   & \dm{Klempner}  & plumber \\
\dm{Elektriker} & electrician & \dm{Mechaniker} & mechanic \\
\dm{Maler}      & painter     & \dm{Fahrer}    & driver \\
\dm{Pilot}      & pilot       & \dm{Soldat}    & soldier \\
\dm{Beamte}     & civil servant & \df{Sekretärin} & secretary \\
\dm{Programmierer} & programmer & \dm{Wissenschaftler} & scientist \\
\dm{Journalist} & journalist  & \dm{Schauspieler} & actor \\
\dm{Musiker}    & musician    & \dm{Künstler}  & artist \\
\dm{Professor}  & professor   & \dm{Direktor}  & director \\
\dm{Mitarbeiter} & employee   & \dm{Kunde}     & customer \\
\df{Firma}      & company     & \dn{Büro}      & office \\
\df{Fabrik}     & factory     & \df{Stelle}    & post, job \\
\dn{Gehalt}     & salary      & \dm{Lohn}      & wages \\
\dm{Vertrag}    & contract    & \df{Bewerbung} & application \\
\dn{Vorstellungsgespräch} & interview & \df{Ausbildung} & training \\
\dm{Urlaub}     & holiday     & \df{Rente}     & pension \\
\df{Besprechung} & meeting    & \df{Arbeitslosigkeit} & unemployment \\
\end{dtbl}

\begin{gnote}
Most job names form the feminine with \textit{-in}:
\textit{Lehrer} $\rightarrow$ \textit{Lehrerin}, plural
\textit{Lehrerinnen}. \textit{der Beamte} is an adjectival noun
(\textit{ein Beamt\textbf{er}}), and \textit{der Polizist}, \textit{der Soldat}
and \textit{der Kunde} are weak nouns.
\end{gnote}

% ------------------------------------------------------------
\gsec{Places in town}

\begin{dtbl}[\footnotesize]{W{1.08}W{0.92}W{1.08}W{0.92}}
\rowcolor{Ink} \hd{Deutsch} & \hd{English} & \hd{Deutsch} & \hd{English} \\ \hdrule
\df{Stadt}      & town, city  & \dn{Dorf}      & village \\
\dn{Zentrum}    & centre      & \df{Straße}    & street \\
\dm{Platz}      & square      & \df{Ecke}      & corner \\
\dn{Gebäude}    & building    & \dn{Rathaus}   & town hall \\
\df{Kirche}     & church      & \dm{Dom}       & cathedral \\
\dn{Schloss}    & palace      & \df{Burg}      & castle \\
\dn{Museum}     & museum      & \dn{Theater}   & theatre \\
\dn{Kino}       & cinema      & \df{Bibliothek} & library \\
\df{Schule}     & school      & \df{Universität} & university \\
\dn{Krankenhaus} & hospital   & \df{Apotheke}  & pharmacy \\
\df{Post}       & post office & \df{Bank}      & bank \\
\dm{Supermarkt} & supermarket & \dm{Markt}     & market \\
\dn{Geschäft}   & shop        & \dm{Laden}     & shop \\
\df{Bäckerei}   & bakery      & \df{Metzgerei} & butcher's \\
\dn{Restaurant} & restaurant  & \dn{Café}      & café \\
\df{Kneipe}     & pub         & \dn{Hotel}     & hotel \\
\dm{Bahnhof}    & station     & \dm{Flughafen} & airport \\
\df{Haltestelle} & bus stop   & \dm{Parkplatz} & car park \\
\dm{Park}       & park        & \dm{Spielplatz} & playground \\
\dn{Schwimmbad} & swimming pool & \dn{Stadion} & stadium \\
\df{Brücke}     & bridge      & \dm{Turm}      & tower \\
\dn{Denkmal}    & monument    & \dm{Friedhof}  & cemetery \\
\df{Polizei}    & police      & \df{Feuerwehr} & fire service \\
\df{Tankstelle} & petrol station & \dm{Zoo}    & zoo \\
\df{Fußgängerzone} & pedestrian zone & \dn{Einkaufszentrum} & shopping centre \\
\df{Wohnung}    & flat        & \dn{Viertel}   & district \\
\df{Umgebung}   & surroundings & \df{Adresse}  & address \\
\df{Grenze}     & border      & \dn{Land}      & country \\
\end{dtbl}

\begin{gnote}
\textit{das Schloss} is a palace or a lock; \textit{die Burg} is a fortified
castle. Both \textit{das Geschäft} and \textit{der Laden} mean \emph{shop}, but
\textit{Geschäft} also means \emph{business}.
\end{gnote}

% ------------------------------------------------------------
\gsec{Travel and transport}

\begin{dtbl}[\footnotesize]{W{1.08}W{0.92}W{1.08}W{0.92}}
\rowcolor{Ink} \hd{Deutsch} & \hd{English} & \hd{Deutsch} & \hd{English} \\ \hdrule
\df{Reise}      & journey     & \dm{Urlaub}    & holiday \\
\dn{Auto}       & car         & \dm{Wagen}     & car, carriage \\
\dm{Bus}        & bus         & \df{Bahn}      & railway \\
\dm{Zug}        & train       & \df{Straßenbahn} & tram \\
\df{U-Bahn}     & underground & \dn{Fahrrad}   & bicycle \\
\dn{Motorrad}   & motorbike   & \dn{Flugzeug}  & aeroplane \\
\dn{Schiff}     & ship        & \df{Fähre}     & ferry \\
\dn{Taxi}       & taxi        & \dm{Lastwagen} & lorry \\
\dm{Bahnhof}    & station     & \dm{Flughafen} & airport \\
\dm{Hafen}      & harbour     & \dn{Gleis}     & platform, track \\
\df{Fahrkarte}  & ticket      & \dm{Fahrplan}  & timetable \\
\df{Abfahrt}    & departure   & \df{Ankunft}   & arrival \\
\df{Verspätung} & delay       & \dm{Anschluss} & connection \\
\dm{Koffer}     & suitcase    & \dn{Gepäck}    & luggage \\
\dm{Reisepass}  & passport    & \dn{Visum}     & visa \\
\df{Karte}      & map, ticket & \dm{Stadtplan} & street map \\
\dn{Hotel}      & hotel       & \df{Jugendherberge} & youth hostel \\
\df{Buchung}    & booking     & \dm{Ausflug}   & excursion \\
\df{Fahrt}      & ride, trip  & \dm{Flug}      & flight \\
\dm{Sitzplatz}  & seat        & \dm{Schaffner} & conductor \\
\df{Autobahn}   & motorway    & \df{Kreuzung}  & junction \\
\df{Ampel}      & traffic light & \dn{Benzin}  & petrol \\
\dm{Führerschein} & driving licence & \dm{Stau} & traffic jam \\
\df{Richtung}   & direction   & \dm{Weg}       & way, path \\
\df{Entfernung} & distance    & \df{Geschwindigkeit} & speed \\
\dm{Unfall}     & accident    & \df{Panne}     & breakdown \\
\dm{Reifen}     & tyre        & \dn{Rad}       & wheel \\
\df{Grenze}     & border      & \dn{Ziel}      & destination \\
\end{dtbl}

\begin{gnote}
\textit{fahren} covers travelling by any vehicle, \textit{fliegen} by air and
\textit{gehen} only on foot. \textit{Ich fahre nach Berlin} says nothing about
whether you drive or take the train.
\end{gnote}

% ------------------------------------------------------------
\gsec{Nature and weather}

\begin{dtbl}[\footnotesize]{W{1.08}W{0.92}W{1.08}W{0.92}}
\rowcolor{Ink} \hd{Deutsch} & \hd{English} & \hd{Deutsch} & \hd{English} \\ \hdrule
\dn{Wetter}     & weather     & \df{Natur}     & nature \\
\df{Sonne}      & sun         & \dm{Mond}      & moon \\
\dm{Stern}      & star        & \dm{Himmel}    & sky \\
\df{Wolke}      & cloud       & \dm{Regen}     & rain \\
\dm{Schnee}     & snow        & \dn{Eis}       & ice \\
\dm{Wind}       & wind        & \dm{Sturm}     & storm \\
\dm{Nebel}      & fog         & \dm{Donner}    & thunder \\
\dm{Blitz}      & lightning   & \dm{Regenbogen} & rainbow \\
\df{Temperatur} & temperature & \dm{Grad}      & degree \\
\df{Hitze}      & heat        & \df{Kälte}     & cold \\
\dm{Frühling}   & spring      & \dm{Sommer}    & summer \\
\dm{Herbst}     & autumn      & \dm{Winter}    & winter \\
\df{Jahreszeit} & season      & \dn{Klima}     & climate \\
\dm{Baum}       & tree        & \dm{Wald}      & forest \\
\df{Blume}      & flower      & \dn{Gras}      & grass \\
\dn{Blatt}      & leaf        & \df{Wurzel}    & root \\
\dm{Berg}       & mountain    & \dn{Tal}       & valley \\
\dm{Hügel}      & hill        & \dm{Fluss}     & river \\
\dm{See}        & lake        & \dn{Meer}      & sea \\
\df{Küste}      & coast       & \dm{Strand}    & beach \\
\df{Insel}      & island      & \df{Wüste}     & desert \\
\df{Wiese}      & meadow      & \dn{Feld}      & field \\
\dm{Stein}      & stone       & \dm{Sand}      & sand \\
\df{Erde}       & earth, soil & \df{Luft}      & air \\
\dn{Feuer}      & fire        & \df{Umwelt}    & environment \\
\df{Pflanze}    & plant       & \dm{Bach}      & stream \\
\dm{Gipfel}     & summit      & \df{Höhle}     & cave \\
\dm{Boden}      & ground      & \df{Landschaft} & landscape \\
\end{dtbl}

\begin{gnote}
Weather is described with impersonal \textit{es}: \textit{es regnet},
\textit{es schneit}, \textit{es ist kalt}. \textit{der See} is a lake but
\textit{die See} is the sea — same word, two genders.
\end{gnote}

% ------------------------------------------------------------
\gsec{Time and the calendar}

\begin{dtbl}[\footnotesize]{W{1.08}W{0.92}W{1.08}W{0.92}}
\rowcolor{Ink} \hd{Deutsch} & \hd{English} & \hd{Deutsch} & \hd{English} \\ \hdrule
\df{Zeit}       & time        & \df{Uhr}       & clock, o'clock \\
\df{Sekunde}    & second      & \df{Minute}    & minute \\
\df{Stunde}     & hour        & \dm{Tag}       & day \\
\df{Woche}      & week        & \dm{Monat}     & month \\
\dn{Jahr}       & year        & \dn{Jahrzehnt} & decade \\
\dn{Jahrhundert} & century    & \dm{Kalender}  & calendar \\
\dm{Morgen}     & morning     & \dm{Vormittag} & late morning \\
\dm{Mittag}     & midday      & \dm{Nachmittag} & afternoon \\
\dm{Abend}      & evening     & \df{Nacht}     & night \\
\df{Mitternacht} & midnight   & \df{Dämmerung} & dusk, dawn \\
\dm{Montag}     & Monday      & \dm{Dienstag}  & Tuesday \\
\dm{Mittwoch}   & Wednesday   & \dm{Donnerstag} & Thursday \\
\dm{Freitag}    & Friday      & \dm{Samstag}   & Saturday \\
\dm{Sonntag}    & Sunday      & \dn{Wochenende} & weekend \\
\dm{Januar}     & January     & \dm{Februar}   & February \\
\dm{März}       & March       & \dm{April}     & April \\
\dm{Mai}        & May         & \dm{Juni}      & June \\
\dm{Juli}       & July        & \dm{August}    & August \\
\dm{September}  & September   & \dm{Oktober}   & October \\
\dm{November}   & November    & \dm{Dezember}  & December \\
\dm{Feiertag}   & public holiday & \dm{Geburtstag} & birthday \\
\dm{Termin}     & appointment & \dm{Augenblick} & moment \\
\df{Vergangenheit} & past     & \df{Gegenwart} & present \\
\df{Zukunft}    & future      & \dm{Anfang}    & beginning \\
\dn{Ende}       & end         & \df{Dauer}     & duration \\
\df{Pause}      & break       & \df{Frist}     & deadline \\
\df{Verspätung} & delay       & \df{Eile}      & hurry \\
\dn{Datum}      & date        & \df{Epoche}    & era \\
\end{dtbl}

\begin{gnote}
Days, months and seasons are all masculine, and all of them take
\textit{am} or \textit{im}: \textit{\textbf{am} Montag}, \textit{\textbf{im}
Mai}, \textit{\textbf{im} Sommer}.
\end{gnote}

% ------------------------------------------------------------
\gsec{School and study}

\begin{dtbl}[\footnotesize]{W{1.08}W{0.92}W{1.08}W{0.92}}
\rowcolor{Ink} \hd{Deutsch} & \hd{English} & \hd{Deutsch} & \hd{English} \\ \hdrule
\df{Schule}     & school      & \df{Universität} & university \\
\df{Klasse}     & class       & \dm{Kurs}      & course \\
\dn{Fach}       & subject     & \dm{Unterricht} & lessons \\
\df{Prüfung}    & exam        & \df{Note}      & mark, grade \\
\dn{Zeugnis}    & report      & \df{Hausaufgabe} & homework \\
\dn{Heft}       & exercise book & \dn{Buch}    & book \\
\dn{Lehrbuch}   & textbook    & \dn{Wörterbuch} & dictionary \\
\dm{Stift}      & pen         & \dm{Bleistift} & pencil \\
\dm{Kugelschreiber} & biro    & \dm{Radiergummi} & rubber \\
\dn{Lineal}     & ruler       & \df{Schere}    & scissors \\
\df{Tafel}      & board       & \df{Kreide}    & chalk \\
\dm{Schreibtisch} & desk      & \dn{Klassenzimmer} & classroom \\
\df{Bibliothek} & library     & \df{Mensa}     & canteen \\
\df{Pause}      & break       & \dpl{Ferien}   & school holidays \\
\dn{Semester}   & semester    & \df{Vorlesung} & lecture \\
\dn{Seminar}    & seminar     & \df{Hausarbeit} & term paper \\
\dn{Studium}    & studies     & \dm{Abschluss} & degree \\
\dm{Schüler}    & pupil       & \df{Schülerin} & pupil \eng{f} \\
\dm{Student}    & student     & \df{Studentin} & student \eng{f} \\
\df{Frage}      & question    & \df{Antwort}   & answer \\
\dm{Fehler}     & mistake     & \dn{Ergebnis}  & result \\
\dn{Wissen}     & knowledge   & \df{Sprache}   & language \\
\df{Mathematik} & maths       & \df{Geschichte} & history \\
\df{Erdkunde}   & geography   & \df{Physik}    & physics \\
\df{Chemie}     & chemistry   & \df{Biologie}  & biology \\
\df{Kunst}      & art         & \df{Musik}     & music \\
\dm{Sport}      & sport       & \df{Informatik} & computing \\
\df{Aufgabe}    & task        & \dn{Beispiel}  & example \\
\end{dtbl}

\begin{gnote}
\textit{die Ferien} exists only in the plural and means school or university
holidays; a personal holiday from work is \textit{der Urlaub}. \textit{der
Student} and \textit{der Schüler} differ: a \textit{Student} is at university.
\end{gnote}

% ------------------------------------------------------------
\gsec{Technology and media}

\begin{dtbl}[\footnotesize]{W{1.08}W{0.92}W{1.08}W{0.92}}
\rowcolor{Ink} \hd{Deutsch} & \hd{English} & \hd{Deutsch} & \hd{English} \\ \hdrule
\dm{Computer}   & computer    & \dm{Laptop}    & laptop \\
\dn{Handy}      & mobile phone & \dn{Smartphone} & smartphone \\
\dm{Bildschirm} & screen      & \df{Tastatur}  & keyboard \\
\df{Maus}       & mouse       & \dm{Drucker}   & printer \\
\df{Datei}      & file        & \dm{Ordner}    & folder \\
\dn{Programm}   & program     & \df{Anwendung} & application \\
\dn{Internet}   & internet    & \df{Webseite}  & website \\
\df{E-Mail}     & email       & \dn{Passwort}  & password \\
\dn{Netzwerk}   & network     & \dm{Akku}      & battery \\
\dn{Ladegerät}  & charger     & \dn{Kabel}     & cable \\
\dm{Lautsprecher} & speaker   & \dm{Kopfhörer} & headphones \\
\df{Kamera}     & camera      & \dn{Foto}      & photo \\
\dn{Video}      & video       & \dm{Fernseher} & television \\
\dn{Radio}      & radio       & \dpl{Nachrichten} & the news \\
\df{Zeitung}    & newspaper   & \df{Zeitschrift} & magazine \\
\dm{Artikel}    & article     & \df{Werbung}   & advertising \\
\dm{Sender}     & channel     & \df{Sendung}   & programme \\
\dm{Film}       & film        & \df{Musik}     & music \\
\dn{Lied}       & song        & \dn{Spiel}     & game \\
\df{Technik}    & technology  & \df{Software}  & software \\
\dpl{Daten}     & data        & \dm{Speicher}  & memory, storage \\
\df{Suche}      & search      & \dm{Anruf}     & phone call \\
\df{Nachricht}  & message     & \dm{Fehler}    & error \\
\dm{Roboter}    & robot       & \df{Erfindung} & invention \\
\dm{Bericht}    & report      & \df{Sicherheit} & security \\
\dm{Nutzer}     & user        & \dn{Konto}     & account \\
\df{Verbindung} & connection  & \df{Einstellung} & setting \\
\dm{Knopf}      & button      & \dm{Schalter}  & switch \\
\end{dtbl}

\begin{gnote}
\textit{das Handy} is a false friend: it means a mobile phone, not something
convenient. \textit{die Nachricht} is a single message; the plural
\textit{die Nachrichten} is the news bulletin.
\end{gnote}

% ------------------------------------------------------------
\gsec{Money and shopping}

\begin{dtbl}[\footnotesize]{W{1.08}W{0.92}W{1.08}W{0.92}}
\rowcolor{Ink} \hd{Deutsch} & \hd{English} & \hd{Deutsch} & \hd{English} \\ \hdrule
\dn{Geld}       & money       & \df{Währung}   & currency \\
\dm{Euro}       & euro        & \dm{Cent}      & cent \\
\df{Münze}      & coin        & \dm{Schein}    & banknote \\
\df{Bank}       & bank        & \dn{Konto}     & account \\
\df{Kreditkarte} & credit card & \dm{Geldautomat} & cash machine \\
\df{Kasse}      & till        & \df{Quittung}  & receipt \\
\dm{Preis}      & price       & \dpl{Kosten}   & costs \\
\df{Rechnung}   & bill        & \dn{Angebot}   & offer \\
\dm{Rabatt}     & discount    & \dm{Ausverkauf} & sale \\
\df{Steuer}     & tax         & \df{Miete}     & rent \\
\dn{Gehalt}     & salary      & \dpl{Schulden} & debts \\
\dm{Kredit}     & loan        & \df{Ersparnis} & saving \\
\dn{Wechselgeld} & change     & \dm{Einkauf}   & shopping \\
\dm{Einkaufswagen} & trolley  & \dm{Korb}      & basket \\
\df{Tüte}       & bag         & \df{Ware}      & goods \\
\dn{Produkt}    & product     & \df{Qualität}  & quality \\
\df{Garantie}   & guarantee   & \dm{Umtausch}  & exchange \\
\df{Lieferung}  & delivery    & \df{Bestellung} & order \\
\df{Bezahlung}  & payment     & \dm{Wert}      & value \\
\dn{Vermögen}   & wealth      & \df{Spende}    & donation \\
\dm{Gewinn}     & profit      & \dm{Verlust}   & loss \\
\dm{Markt}      & market      & \dn{Schaufenster} & shop window \\
\df{Größe}      & size        & \df{Auswahl}   & selection \\
\dm{Verkäufer}  & seller      & \dm{Käufer}    & buyer \\
\df{Versicherung} & insurance & \df{Gebühr}    & fee \\
\dn{Bargeld}    & cash        & \df{Überweisung} & bank transfer \\
\dm{Beleg}      & receipt     & \df{Werbung}   & advertising \\
\df{Wirtschaft} & economy     & \dm{Handel}    & trade \\
\end{dtbl}

\begin{gnote}
\textit{die Kosten}, \textit{die Schulden} and \textit{die Zinsen} exist only in
the plural. Note that \textit{Geld} is uncountable: \textit{Ich habe kein Geld},
never \textit{keine Gelder} in ordinary use.
\end{gnote}

% ------------------------------------------------------------
\gsec{Health and illness}

\begin{dtbl}[\footnotesize]{W{1.08}W{0.92}W{1.08}W{0.92}}
\rowcolor{Ink} \hd{Deutsch} & \hd{English} & \hd{Deutsch} & \hd{English} \\ \hdrule
\df{Gesundheit} & health      & \df{Krankheit} & illness \\
\dm{Arzt}       & doctor      & \df{Praxis}    & surgery \\
\dn{Krankenhaus} & hospital   & \df{Apotheke}  & pharmacy \\
\dm{Termin}     & appointment & \df{Untersuchung} & examination \\
\df{Behandlung} & treatment   & \df{Diagnose}  & diagnosis \\
\dn{Rezept}     & prescription & \dn{Medikament} & medicine \\
\df{Tablette}   & tablet      & \df{Salbe}     & ointment \\
\df{Spritze}    & injection   & \dm{Verband}   & bandage \\
\dn{Pflaster}   & plaster     & \df{Operation} & operation \\
\dm{Schmerz}    & pain        & \dn{Fieber}    & fever \\
\df{Erkältung}  & cold        & \df{Grippe}    & flu \\
\dm{Husten}     & cough       & \dm{Schnupfen} & runny nose \\
\dpl{Halsschmerzen} & sore throat & \dpl{Kopfschmerzen} & headache \\
\dpl{Bauchschmerzen} & stomach ache & \dpl{Zahnschmerzen} & toothache \\
\df{Allergie}   & allergy     & \df{Verletzung} & injury \\
\dm{Unfall}     & accident    & \dm{Notfall}   & emergency \\
\dm{Krankenwagen} & ambulance & \df{Wunde}     & wound \\
\dm{Bruch}      & fracture    & \df{Narbe}     & scar \\
\df{Impfung}    & vaccination & \dn{Blut}      & blood \\
\dm{Puls}       & pulse       & \dm{Blutdruck} & blood pressure \\
\df{Ernährung}  & diet        & \df{Bewegung}  & exercise \\
\dm{Stress}     & stress      & \df{Erholung}  & recovery \\
\dm{Schlaf}     & sleep       & \df{Müdigkeit} & tiredness \\
\df{Übelkeit}   & nausea      & \dm{Durchfall} & diarrhoea \\
\dm{Zahnarzt}   & dentist     & \df{Vorsorge}  & prevention \\
\df{Krankenkasse} & health insurance & \df{Pflege} & care \\
\dm{Patient}    & patient     & \df{Station}   & ward \\
\df{Genesung}   & recovery    & \df{Therapie}  & therapy \\
\end{dtbl}

\begin{gnote}
Aches are usually plural: \textit{Ich habe Kopfschmerzen}. To say what hurts,
German uses \textit{wehtun} with the dative: \textit{Mir tut der Kopf weh}.
\end{gnote}

% ------------------------------------------------------------
\gsec{Countries and languages}

\begin{dtbl}[\footnotesize]{W{1.02}W{0.98}W{1.02}W{0.98}}
\rowcolor{Ink} \hd{Land} & \hd{Country} & \hd{Sprache} & \hd{Language} \\ \hdrule
Deutschland     & Germany     & Deutsch        & German \\
Österreich      & Austria     & Deutsch        & German \\
\df{Schweiz}    & Switzerland & Deutsch        & German \\
England         & England     & Englisch       & English \\
Schottland      & Scotland    & Englisch       & English \\
Irland          & Ireland     & Englisch       & English \\
Frankreich      & France      & Französisch    & French \\
Italien         & Italy       & Italienisch    & Italian \\
Spanien         & Spain       & Spanisch       & Spanish \\
Portugal        & Portugal    & Portugiesisch  & Portuguese \\
\dpl{Niederlande} & Netherlands & Niederländisch & Dutch \\
Belgien         & Belgium     & Flämisch       & Flemish \\
Dänemark        & Denmark     & Dänisch        & Danish \\
Schweden        & Sweden      & Schwedisch     & Swedish \\
Norwegen        & Norway      & Norwegisch     & Norwegian \\
Finnland        & Finland     & Finnisch       & Finnish \\
Polen           & Poland      & Polnisch       & Polish \\
Tschechien      & Czechia     & Tschechisch    & Czech \\
Ungarn          & Hungary     & Ungarisch      & Hungarian \\
Griechenland    & Greece      & Griechisch     & Greek \\
\df{Türkei}     & Turkey      & Türkisch       & Turkish \\
Russland        & Russia      & Russisch       & Russian \\
\df{Ukraine}    & Ukraine     & Ukrainisch     & Ukrainian \\
China           & China       & Chinesisch     & Chinese \\
Japan           & Japan       & Japanisch      & Japanese \\
Indien          & India       & Hindi          & Hindi \\
Ägypten         & Egypt       & Arabisch       & Arabic \\
Brasilien       & Brazil      & Portugiesisch  & Portuguese \\
\dpl{USA}       & USA         & Englisch       & English \\
\end{dtbl}

\begin{gnote}
Most country names take no article. The handful that do are shown with it here,
and they keep it in use: \textit{in \textbf{der} Schweiz}, \textit{in
\textbf{die} Türkei}. Language names are capitalised nouns; the matching
adjectives (\textit{deutsch}, \textit{französisch}) are not.
\end{gnote}

% ------------------------------------------------------------
\gsec{Colours, shapes and materials}

\tcap{Colours \eng{adjectives, so no article}}

\begin{dtbl}[\footnotesize]{W{0.85}W{1.15}W{0.85}W{1.15}}
\rowcolor{Ink} \hd{Deutsch} & \hd{English} & \hd{Deutsch} & \hd{English} \\ \hdrule
rot        & red         & blau      & blue \\
grün       & green       & gelb      & yellow \\
schwarz    & black       & weiß      & white \\
grau       & grey        & braun     & brown \\
orange     & orange      & rosa      & pink \\
lila       & purple      & türkis    & turquoise \\
beige      & beige       & bunt      & colourful \\
hell       & light       & dunkel    & dark \\
golden     & golden      & silbern   & silver \\
\end{dtbl}

\tcap{Shapes and materials}

\begin{dtbl}[\footnotesize]{W{1.08}W{0.92}W{1.08}W{0.92}}
\rowcolor{Ink} \hd{Deutsch} & \hd{English} & \hd{Deutsch} & \hd{English} \\ \hdrule
\df{Form}       & shape       & \dm{Kreis}     & circle \\
\dn{Quadrat}    & square      & \dn{Rechteck}  & rectangle \\
\dn{Dreieck}    & triangle    & \df{Linie}     & line \\
\dm{Punkt}      & point, dot  & \df{Kugel}     & sphere \\
\dm{Würfel}     & cube        & \dm{Zylinder}  & cylinder \\
\dn{Holz}       & wood        & \dn{Metall}    & metal \\
\dn{Eisen}      & iron        & \dm{Stahl}     & steel \\
\dn{Gold}       & gold        & \dn{Silber}    & silver \\
\dn{Glas}       & glass       & \dn{Plastik}   & plastic \\
\dn{Papier}     & paper       & \dm{Stein}     & stone \\
\dm{Beton}      & concrete    & \dn{Gummi}     & rubber \\
\end{dtbl}

\begin{gnote}
\textit{rosa}, \textit{lila}, \textit{orange} and \textit{beige} never take
adjective endings: \textit{ein \textbf{rosa} Kleid}, not \textit{ein rosanes
Kleid}. Every other colour declines normally.
\end{gnote}

% ------------------------------------------------------------
\gsec{Common adjectives}

\gintro{Listed in opposite pairs, which is how they are most easily remembered.
None of these takes an ending in the form shown; endings appear only in front of
a noun.}

\begin{dtbl}[\footnotesize]{W{0.95}W{1.05}W{0.95}W{1.05}}
\rowcolor{Ink} \hd{Deutsch} & \hd{English} & \hd{Gegenteil} & \hd{English} \\ \hdrule
groß        & big          & klein        & small \\
lang        & long         & kurz         & short \\
hoch        & high         & niedrig      & low \\
breit       & wide         & schmal       & narrow \\
dick        & thick, fat   & dünn         & thin \\
schwer      & heavy, hard  & leicht       & light, easy \\
alt         & old          & jung         & young \\
neu         & new          & gebraucht    & used \\
gut         & good         & schlecht     & bad \\
richtig     & right        & falsch       & wrong \\
schön       & beautiful    & hässlich     & ugly \\
sauber      & clean        & schmutzig    & dirty \\
warm        & warm         & kalt         & cold \\
heiß        & hot          & kühl         & cool \\
hell        & bright       & dunkel       & dark \\
laut        & loud         & leise        & quiet \\
schnell     & fast         & langsam      & slow \\
früh        & early        & spät         & late \\
teuer       & expensive    & billig       & cheap \\
reich       & rich         & arm          & poor \\
stark       & strong       & schwach      & weak \\
hart        & hard         & weich        & soft \\
voll        & full         & leer         & empty \\
offen       & open         & geschlossen  & closed \\
einfach     & easy, simple & schwierig    & difficult \\
interessant & interesting  & langweilig   & boring \\
glücklich   & happy        & traurig      & sad \\
freundlich  & friendly     & unfreundlich & unfriendly \\
\end{dtbl}

\begin{gnote}
\textit{schwer} and \textit{leicht} each cover two English words:
\emph{heavy/difficult} and \emph{light/easy}. \textit{hoch} loses its
\textit{c} before an ending: \textit{ein \textbf{hoher} Berg}.
\end{gnote}

% ------------------------------------------------------------
\gsec{Everyday verbs}

\gintro{Regular verbs are set plain; \vi{irregular verbs} are bold and
coloured. An irregular verb changes its stem in the past, and often in the
\textit{du} and \textit{er} forms of the present as well. Their full forms are
in Part XV.}

\begin{dtbl}[\footnotesize]{W{1.02}W{0.98}W{1.02}W{0.98}}
\rowcolor{Ink} \hd{Deutsch} & \hd{English} & \hd{Deutsch} & \hd{English} \\ \hdrule
\vi{sein}       & to be        & \vi{haben}     & to have \\
\vi{werden}     & to become    & \vr{machen}    & to do, make \\
\vi{gehen}      & to go        & \vi{kommen}    & to come \\
\vi{fahren}     & to drive, go & \vi{laufen}    & to run, walk \\
\vi{essen}      & to eat       & \vi{trinken}   & to drink \\
\vi{sehen}      & to see       & \vi{lesen}     & to read \\
\vi{schreiben}  & to write     & \vi{sprechen}  & to speak \\
\vr{sagen}      & to say       & \vr{fragen}    & to ask \\
\vr{antworten}  & to answer    & \vi{nehmen}    & to take \\
\vi{geben}      & to give      & \vi{finden}    & to find \\
\vi{bleiben}    & to stay      & \vi{schlafen}  & to sleep \\
\vi{helfen}     & to help      & \vi{tragen}    & to carry, wear \\
\vi{treffen}    & to meet      & \vi{denken}    & to think \\
\vi{bringen}    & to bring     & \vi{wissen}    & to know \eng{facts} \\
\vi{kennen}     & to know \eng{people} & \vi{stehen} & to stand \\
\vi{sitzen}     & to sit       & \vi{liegen}    & to lie \\
\vr{arbeiten}   & to work      & \vr{lernen}    & to learn \\
\vr{spielen}    & to play      & \vr{kaufen}    & to buy \\
\vr{bezahlen}   & to pay       & \vr{kosten}    & to cost \\
\vr{wohnen}     & to live, reside & \vr{leben}  & to live \\
\vr{lieben}     & to love      & \vr{brauchen}  & to need \\
\vr{suchen}     & to look for  & \vr{zeigen}    & to show \\
\vr{hören}      & to hear      & \vr{warten}    & to wait \\
\vr{kochen}     & to cook      & \vr{lachen}    & to laugh \\
\vr{weinen}     & to cry       & \vr{tanzen}    & to dance \\
\vr{besuchen}   & to visit     & \vr{bestellen} & to order \\
\vr{danken}     & to thank     & \vr{glauben}   & to believe \\
\vi{tun}        & to do        & \vi{lassen}    & to let, leave \\
\end{dtbl}

\begin{gnote}
\textit{wissen} takes a fact or a clause (\textit{Ich weiß, dass\dots}), while
\textit{kennen} takes a person or place you are acquainted with. \textit{leben}
is to be alive; \textit{wohnen} is to reside somewhere.
\end{gnote}

% ------------------------------------------------------------
\gsec{Verbs of movement, speech and thought}

\begin{dtbl}[\footnotesize]{W{1.02}W{0.98}W{1.02}W{0.98}}
\rowcolor{Ink} \hd{Deutsch} & \hd{English} & \hd{Deutsch} & \hd{English} \\ \hdrule
\vi{fliegen}    & to fly       & \vi{schwimmen} & to swim \\
\vi{springen}   & to jump      & \vi{steigen}   & to climb \\
\vi{fallen}     & to fall      & \vi{ziehen}    & to pull \\
\vr{drücken}    & to push      & \vi{werfen}    & to throw \\
\vi{fangen}     & to catch     & \vr{holen}     & to fetch \\
\vr{legen}      & to lay       & \vr{stellen}   & to place \\
\vr{setzen}     & to set, put  & \vi{hängen}    & to hang \\
\vr{öffnen}     & to open      & \vi{schließen} & to close \\
\vi{beginnen}   & to begin     & \vr{enden}     & to end \\
\vr{aufhören}   & to stop      & \vi{anfangen}  & to start \\
\vi{rufen}      & to call      & \vi{schreien}  & to shout \\
\vi{singen}     & to sing      & \vr{flüstern}  & to whisper \\
\vr{erzählen}   & to tell      & \vr{erklären}  & to explain \\
\vr{diskutieren} & to discuss  & \vr{übersetzen} & to translate \\
\vi{verstehen}  & to understand & \vi{vergessen} & to forget \\
\vr{erinnern}   & to remind    & \vi{entscheiden} & to decide \\
\vi{empfehlen}  & to recommend & \vr{hoffen}    & to hope \\
\vr{wünschen}   & to wish      & \vr{meinen}    & to think, mean \\
\vi{gefallen}   & to please    & \vr{gehören}   & to belong to \\
\vi{verlieren}  & to lose      & \vi{gewinnen}  & to win \\
\vi{vergleichen} & to compare  & \vr{ändern}    & to change \\
\vr{versuchen}  & to try       & \vr{üben}      & to practise \\
\vr{benutzen}   & to use       & \vr{prüfen}    & to check \\
\vi{halten}     & to hold, stop & \vi{waschen}  & to wash \\
\vi{tragen}     & to carry     & \vi{wachsen}   & to grow \\
\vi{sterben}    & to die       & \vi{schneiden} & to cut \\
\vr{reisen}     & to travel    & \vr{packen}    & to pack \\
\vi{schlagen}   & to hit       & \vr{feiern}    & to celebrate \\
\end{dtbl}

\begin{gnote}
\textit{legen}, \textit{stellen} and \textit{setzen} are regular and take an
accusative object; their irregular partners \textit{liegen}, \textit{stehen} and
\textit{sitzen} describe the resulting position and take the dative.
\end{gnote}
